\documentclass[a4paper,12pt,oneside,onecolumn,fleqn]{report}  % Sayfa biçimi ve metin puntosu
\usepackage{ucs}
\usepackage[utf8]{inputenc}
%\usepackage[style=ieee,sorting=none]{biblatex}%%Added by RC-FUN BABEL ÜZERİNDE OLACAK
\usepackage[english,turkish]{babel} 

\usepackage[square,numbers]{natbib}
% \bibliographystyle{unsrtnat} %{apalike} 
\usepackage{notoccite}

\usepackage[version=4]{mhchem}

\usepackage[absolute,overlay]{textpos}
\usepackage[document]{ragged2e}	%Sağa sola ve heriki yana yasla 
\usepackage[dvipsnames]{xcolor}	%renk paketi
\usepackage[top=3.0cm, bottom=2.5cm, left=3.25cm, right=2.5cm]{geometry} % Kenar boşlukları 
\usepackage{textgreek} % Yunanca karakter 
\usepackage{titlesec} % Anabölüm ve alt bölüm başlıkları font ve boşluk değişimi
\usepackage{pdflscape}

\usepackage{graphicx} %[dvips][pdftex]
\usepackage{epstopdf}
% \DeclareGraphicsRule{.pdf}{pdf}{-crop.pdf}{`pdfcrop #1}
% \usepackage[export]{adjustbox}
\graphicspath{{./images/}}%{{./images/}} %windows için graphicspath{ {c:/.../images/} }
\PassOptionsToPackage{cmyk}{xcolor}

\usepackage{amsmath} 		% principal package in the AMS-LATEX distribution packages amsbsy (for bold symbols), amsopn (for operator names) and amstext (for text embedded in mathematics) are also loaded.
\usepackage{amsfonts}		% extended set of fonts for use in mathematics, including: extra mathematical symbols; blackboard bold letters (uppercase only)
\usepackage{amssymb}		% additional math symbols, like arrows, operators, special characters, geometric figures
\usepackage{amsthm}			% theorem setup typically needed in American Mathematical Society publications
\usepackage{mathtools}		% useful tools for mathematical typesetting, based on amsmath and fixes various deficiencies of amsmath and standard LATEX
\usepackage{mathptmx}		% use Times as default text font, and provide maths support
\usepackage{bm}				% defines a command \bm which makes its argument bold, the argument may be any maths object from a single symbol to an expression
\usepackage{lmodern}		% the Latin Modern family of fonts consists of 72 text fonts and 20 mathematics fonts, and is based on the Computer Modern fonts released into public domain by AMS
\usepackage{textcomp}		% supports the Text Companion fonts, which provide many text symbols (such as baht, bullet, copyright, musicalnote, onequarter, section, and yen), in the TS1 encoding
\usepackage{setspace}		% set space between lines
\usepackage{indentfirst}	% indent first paragraph after section header
% \usepackage{calculator}		% make calculations inside latex source code +++++++++++++++
\usepackage{textcmds}		% provides shorthand commands for all the text symbols that are traditionally produced in LATEX documents by non-letter ligatures
% \usepackage{subfigure} 
\usepackage{url}			% a form of verbatim command that allows linebreaks at certain characters or combinations of characters, accepts reconfiguration, and can usually be used in the argument to another command
\usepackage{fancyvrb}		% inline verbatim
\usepackage{nameref}		% isimle referans vermek için \nameref{}
\usepackage{mathrsfs}
\usepackage[framed]{matlab-prettifier} % .m file kodları yazmak için
\usepackage{longtable}		%  uzun uzun donanım kısaltmalar listesi için
\usepackage{lipsum}      %  dummy text üreteci
%%%%%%%FU
\usepackage[figurewithin=none]{caption} % Şekiller listesinde yeni bölüm başladığındaki ara boşluğu siler.
\usepackage[tablewithin=none]{caption}  % Tablolar listesinde yeni bölüm başladığındaki ara boşluğu siler.
%%%%%%%%FU
\usepackage{tocloft} %	Tablo,Şekil, İçindekiler Sayfası ayarları
\usepackage[labelsep=quad]{caption}		% provides many ways to customise the captions in floating environments like figure and table, and cooperates with many other packages
\usepackage{subcaption}		% provides a means of using facilities analogous to those of the caption package, when writing captions for subfigures and the like
\usepackage{float} 
% \newcommand*{\nextfloatpart}{\vskip{5em}}
\usepackage{booktabs} %	Tablo alt çizgi kalınlığı ayarı
\usepackage{pgfplots}
\pgfplotsset{compat=1.8}
\usepackage{tikz} %	Tik sembolü ekleme
\usepackage[siunitx, european resistor, european voltage, european current]{circuitikz}
\usetikzlibrary{arrows,shapes,positioning}
\usetikzlibrary{decorations.markings,decorations.pathmorphing,decorations.pathreplacing}
\usetikzlibrary{calc,patterns,shapes.geometric,babel} 
\pdfpageattr {/Group << /S /Transparency /I true /CS /DeviceRGB>>} 
\usepackage{enumitem}
\usepackage{appendix}
% \usepackage[style=ieee,sorting=none]{biblatex} % fun kaynakça format değişikliği için gerekli
% \addbibresource{./bib/thesis.bib} 

\usepackage{esdiff}			% typesetting of derivatives


\renewcommand{\theequation}{\arabic{chapter}. \arabic{equation}} % fun
\usepackage{adjustbox}		% provides several macros to adjust boxed content

\usepackage{minted}			% expressive syntax highlighting in LATEX using the powerful Pygments library. The package also provides options to customize the highlighted source code output using fancyvrb, requires shell escape
\usepackage{algorithmic}	% older algorithm package
\usepackage{algorithm}		% defines a floating algorithm environment designed to work with the algorithmic style
\usepackage{multicol}		% defines a multicols environment which typesets text in multiple columns
\usepackage{multirow}		% tabular cells spanning multiple rows
\usepackage[flushleft]{threeparttable} % fun alt açıklama notu olan tablo için gerekli
\usepackage{hhline} % fun tabloda ara satırlarda çift çizgi için gerekli
\usepackage{array} % fun tabloda ara satırlarda çift çizgi için gerekli

% \usepackage{showframe} %bölgeleri çizerek belirtir. +++++++++++++++++++++++++++ \theta


\newcolumntype{L}[1]{>{\raggedright\let\newline\\\arraybackslash\hspace{0pt}}m{#1}} % fun tabloda ara satırlarda çift çizgi için gerekli
\newcolumntype{C}[1]{>{\centering\let\newline\\\arraybackslash\hspace{0pt}}m{#1}} % fun tabloda ara satırlarda çift çizgi için gerekli
\newcolumntype{R}[1]{>{\raggedleft\let\newline\\\arraybackslash\hspace{0pt}}m{#1}} % fun tabloda ara satırlarda çift çizgi için gerekli
%----------------------------------------------------------------------------------------------------
%	Alttaki blok kalın büyük-küçük (bold small caps) fontu 
%	oluşturmayı amaçlar 
\RequirePackage{fix-cm}
\usepackage[T1]{fontenc}
\usepackage{lmodern}
\rmfamily
\DeclareFontShape{T1}{lmr}{b}{sc}{<->ssub*cmr/bx/sc}{}
\DeclareFontShape{T1}{lmr}{bx}{sc}{<->ssub*cmr/bx/sc}{}
%----------------------------------------------------------------------------------------------------
%----------------------------------------------------------------------------------------------------

%----------------------------------------------------------------------------------------------------
%	Alttaki blok matematik denklemlerinin 1cm girdi biçimde sola hizlanmasını amaçlar 
\makeatletter
\newcommand{\mathleft}{\@fleqntrue\@mathmargin1cm}
\newcommand{\mathcenter}{\@fleqnfalse}
% \patchcmd{\table}{10\p@}{1\p@}{}{}% \belowcaptionskip @ 10pt -> 5pt
% \patchcmd{\figure}{10\p@}{1\p@}{}{}% \belowcaptionskip @ 10pt -> 5pt
\makeatother
%----------------------------------------------------------------------------------------------------
%----------------------------------------------------------------------------------------------------

\urlstyle{rm} %	url referans fontu değişimi (şaun kullanılan font roman )

\usepackage{longtable} %	Uzun tabloların bir sonraki sayfadan devam etmesini sağlar

%----------------------------------------------------------------------------------------------------
%	Alttaki blok metni belirli boyutlara sahip bir kutu içerisine sıkıştırmayı amaçlar 
\usepackage{environ}
\usepackage{relsize}
\makeatletter
\NewEnviron{fitbox}[2]{%
	\minipage{#1}%
	\sbox0{\minipage{#1}\strut\BODY\strut\endminipage}%
	\Gscale@div\factor{#2}{\dimexpr\ht0+\dp0\relax}%
	\relscale{\factor}%
	\BODY
	\endminipage
}
\makeatother
%----------------------------------------------------------------------------------------------------
%----------------------------------------------------------------------------------------------------

%----------------------------------------------------------------------------------------------------
% 	Sayfaya sığmayan kelimelerin "-" ile bölünmesini engelleyen kod bloğu
\tolerance=1
\emergencystretch=\maxdimen
\hyphenpenalty=10000
\hbadness=10000
%----------------------------------------------------------------------------------------------------
%----------------------------------------------------------------------------------------------------
\usepackage{csquotes}		% advanced facilities for inline and display quotations
%----------------------------------------------------------------------------------------------------
%	Sayfa numarası boşluğu
\usepackage{calc}
\setlength{\footskip}{\paperheight
	-(1in+\voffset+\topmargin+\headheight+\headsep+\textheight)
	-1.5cm}
%----------------------------------------------------------------------------------------------------+++++++++++++++++++++++++++++++++++

%You can modify the following lengths, which affect all floats.;

% \floatsep: space left between floats (12.0pt plus 2.0pt minus 2.0pt).
% \textfloatsep: space between last top float or first bottom float and the text (20.0pt plus 2.0pt minus 4.0pt).
% \intextsep : space left on top and bottom of an in-text float (12.0pt plus 2.0pt minus 2.0pt).
% \dbltextfloatsep is \textfloatsep for 2 column output (20.0pt plus 2.0pt minus 4.0pt).
% \dblfloatsep is \floatsep for 2 column output (12.0pt plus 2.0pt minus 2.0pt).
% \abovecaptionskip: space above caption (10.0pt).
% \belowcaptionskip: space below caption (0.0pt).

%----------------------------------------------------------------------------------------------------+++++++++++++++++++++++++++++++++++

% \usepackage{etoolbox}
% \makeatletter
% \patchcmd{\thebibliography}{\c@NAT@ctr\z@}{\c@NAT@ctr\z@
%   \renewcommand{\makelabel}[1]{\makebox[1.2em][r]{##1\hfill}}% 
%   \setlength{\labelwidth}{1.2em}%
%   \setlength{\labelsep}{.5em}%
%   \setlength{\leftmargin}{\dimexpr\labelwidth+\labelsep}%
% }{}{}
% \makeatother

%----------------------------------------------------------------------------------------------------
%	Sayfa numarası font değişimi
\usepackage{fancyhdr}
\fancypagestyle{plain}{%
	\fancyhf{} 
	\fancyfoot[C]{\fontsize{10pt}{10pt}\selectfont\thepage} 
	\renewcommand{\headrulewidth}{0pt}
	\renewcommand{\footrulewidth}{0pt}}
\pagestyle{plain}
%----------------------------------------------------------------------------------------------------
\usepackage{reftr}
\usepackage{hyperref}		% to handle cross-referencing commands in LATEX to produce hypertext links in the document
\usepackage{cleveref}		% enhances LATEX’s cross-referencing features, allowing the format of references to be determined automatically according to the type of reference
\crefformat{equation}{Denklem~#2#1#3} %<--- will give Eq. (2.1) 
% \crefformat{equation}{Denklem~(#2#1#3)} %<--- will give Eq. (2.1) 

%----------------------------------------------------------------------------------------------------
%\setlength{\cftsubsecnumwidth}{0.0cm}
%\setlength{\cftsubsecindent}{0.60cm}
%\renewcommand{cftsubsecleader}{}
%-----------------------
%%%%%FU
\usepackage{chngcntr} % Şekil ve tablo numaralandırmasını bölüme başlarında yeni bölüm için sıfırlama işlemi için
\counterwithin{figure}{chapter} % Şekil numaralandırmasını bölüme başlarında yeni bölüm için sıfırlama işlemi için
\counterwithin{table}{chapter} % Tablo numaralandırmasını bölüme başlarında yeni bölüm için sıfırlama işlemi için
%%%%%FU

\begin{document}


%----------------------------------------------------------------------------------------------------
%	Büyük küçük yazı sitilinde "i" karakteri sorununun çözümü için
%	gerekli kod bloğu. Metin içerisinde büyük-küçük yazı sitilinde
%	"i" karekteri kullanmak için "\.\i" yazmanız gerekmektedir.
\expandafter\let\csname\string\T1\string\.-\string\i\endcsname\relax
\shorthandoff{=}

%----------------------------------------------------------------------------------------------------
%----------------------------------------------------------------------------------------------------
% 	Latex dökümanınızı oluşturmadan önce ilk olarak "tezbilgileri.tex" dosyasına girerek
%						kendi metninize ait bilgileri giriniz! 
%	
%	"tezbilgileri.tex"  doyasındaki bilgilerin doğru bir biçimde girilmesi; pencere altı kapak, 
%	iç kapak, onay metni ve beyan metninin otomatik olarak doldurulmasını sağlayacaktır.
%	Herhangi bir problem yaşamadığınız takdirde "pencerealtikapak.tex", "ickapak.tex", "onay.tex"
%	"beyan.tex" dosyalarında herhangi bir değişiklik yapmayınız.

\input{data.tex}
\input{./front/hlf_title.tex}
\input{./front/title.tex}
\input{./front/apprvl.tex}
\input{./front/dclrtn.tex} 
%----------------------------------------------------------------------------------------------------
%----------------------------------------------------------------------------------------------------



%----------------------------------------------------------------------------------------------------
%----------------------------------------------------------------------------------------------------
%			 İçindekiler sayfası ayarları
%	(Gerek duyulmadıkça ayarlarda değişiklik yapmayınız!)
%----------------------------------------------------------------------------------------------------
%----------------------------------------------------------------------------------------------------
\renewcommand{\contentsname}{\hfill {\Large \bfseries\scshape İç\.\i ndek\.\i ler} \hfill}  
\renewcommand{\cfttoctitlefont}{\footnotesize}
\renewcommand{\cftchapdotsep}{\cftdotsep}

\renewcommand{\cftaftertoctitle} { \vspace{-1.25cm}} 
\renewcommand{\cftbeforetoctitleskip}{-1cm}
\addtocontents{toc}{~\hfill\footnotesize\textbf{Sayfa}\vspace{-0.3cm}} 
\addtocontents{toc}{~\hfill\rule{\textwidth}{0.3pt}\par}	%	Sayfa numarası üstüne çizgi ekleme
\renewcommand{\cftchapfont}{\scshape\bfseries\small}%	Ana bölüm başlığı fontu değişimi
\renewcommand{\cftsecfont}{\footnotesize}	%	2. derece alt başlık fontu değişimi
\renewcommand{\cftsubsecfont}{\footnotesize}	%	3. derece alt başlık fontu değişimi
\renewcommand\thechapter{\arabic{chapter}.}
\renewcommand\thesection{\thechapter\arabic{section}.}
\renewcommand\thesubsection{\thesection\arabic{subsection}.}
\setlength{\cftbeforechapskip}{3pt} %	TOC daki 	ana bölüm /chapters'lardan sonra gellen boşluk ayarı
\setlength{\cftchapnumwidth}{0.61cm} %	TOC daki 	ana bölüm  numarası ile chapter başlığı arasındaki genişlik ayarı
\setlength{\cftsecnumwidth}{0.80cm}	%	TOC daki 2. derece alt başlık numarası ile ana bölüm başlığı arasındaki genişlik ayarı
\setlength{\cftsubsecnumwidth}{0.95cm}
\setlength{\cftsubsecindent}{1.45cm}

%----------------------------------------------------------------------------------------------------
%----------------------------------------------------------------------------------------------------
%			Şekiller listesi sayfa ayarları
%	(Gerek duyulmadıkça ayarlarda değişiklik yapmayınız!)
%----------------------------------------------------------------------------------------------------
%----------------------------------------------------------------------------------------------------

\renewcommand\listfigurename{\hfill {\Large \bfseries\scshape Şek\.\i ller L\.\i stes\.\i } \hfill}
\renewcommand{\cftloftitlefont}{\footnotesize}
\renewcommand{\cftchapdotsep}{\cftdotsep}
\renewcommand{\cftafterloftitle}{\vspace{-1cm}} 
\renewcommand{\cftbeforeloftitleskip}{-1cm}
\addtocontents{lof}{~\hfill\footnotesize\textbf{Sayfa}\vspace{-0.4cm}} 
\addtocontents{lof}{~\hfill\rule{\textwidth}{0.3pt}\par\vspace{-0.10cm}}  % Sayfa numarası üstüne çizgi ekleme
\renewcommand{\cftfigfont}{\footnotesize} % Şekil listesi yazı sitili değişimi
\renewcommand\thefigure{\thechapter\arabic{figure}.} 
\renewcommand{\cftfigpresnum}{\bfseries Şekil\:}
\setlength{\cftfignumwidth}{2cm}   
\setlength{\cftfigindent}{0cm}  

%----------------------------------------------------------------------------------------------------
%----------------------------------------------------------------------------------------------------
%				Tablo listesi sayfa ayarları
%	(Gerek duyulmadıkça ayarlarda değişiklik yapmayınız!)
%----------------------------------------------------------------------------------------------------
%----------------------------------------------------------------------------------------------------

\renewcommand\listtablename{\hfill {\Large \bfseries\scshape Tablolar L\.\i stes\.\i} \hfill}
\renewcommand{\cftlottitlefont}{\footnotesize}
\renewcommand{\cftchapdotsep}{\cftdotsep}
\renewcommand{\cftafterlottitle}{\vspace{-1cm}} 
\renewcommand{\cftbeforelottitleskip}{-1cm}
\addtocontents{lot}{~\hfill\footnotesize\textbf{Sayfa}\vspace{-0.4cm}} 
\addtocontents{lot}{~\hfill\rule{\textwidth}{0.3pt}\par\vspace{-0.10cm}}  % Sayfa numarası üstüne çizgi ekleme
\renewcommand{\cfttabfont}{\footnotesize} %	Tablo listesi yazı sitili değişimi
\renewcommand\thetable{\thechapter\arabic{table}.} 
\renewcommand{\cfttabpresnum}{\bfseries Tablo\:}
\setlength{\cfttabnumwidth}{2cm}  
\setlength{\cfttabindent}{0cm}  
\setlength{\cftbeforetabskip}{3pt}

\makeatletter %fun
\def\@seccntformat#1{\@ifundefined{#1@cntformat}%
   {\csname the#1\endcsname\space}%    default
   {\csname #1@cntformat\endcsname}}%  enable individual control
\newcommand\section@cntformat{\thesection.\space}       % section-level
\newcommand\subsection@cntformat{\thesubsection.\space} % subsection-level
\makeatother



%----------------------------------------------------------------------------------------------------
%----------------------------------------------------------------------------------------------------
%			Şekil ve Tablo yazıları format ve stil ayarları
%		(Gerek duyulmadıkça ayarlarda değişiklik yapmayınız!)
%----------------------------------------------------------------------------------------------------
%----------------------------------------------------------------------------------------------------

\captionsetup{format=hang, font=scriptsize, labelfont=bf, labelsep=space,aboveskip=4mm, belowskip=8mm} % caption tools
\renewcommand{\tablename}{Tablo} % change table to Tablo
\renewcommand{\figurename}{Şekil}  % change Figure to Şekil
%----------------------------------------------------------------------------------------------------
%----------------------------------------------------------------------------------------------------


%----------------------------------------------------------------------------------------------------
%----------------------------------------------------------------------------------------------------
%			Giriş sayfasına kadar olan Bölümbaşlıkları format ve stil ayarları 
%----------------------------------------------------------------------------------------------------
%----------------------------------------------------------------------------------------------------

\titleformat{\chapter}[hang]{\centering\Large\bfseries\scshape}{\thechapter}{0.01cm}{\Large\bfseries\scshape}
\titlespacing{\chapter}{0pt}{-1cm}{20pt}
\titlespacing*{\chapter}{10pt}{-1cm}{20pt}
% \titlelabel{\thetitle. } 
%----------------------------------------------------------------------------------------------------
%----------------------------------------------------------------------------------------------------

\pagenumbering{roman}
\setcounter{page}{4}
%----------------------------------------------------------------------------------------------------
%----------------------------------------------------------------------------------------------------
%		Lütfen "onsoz.text" dosyasına girerek Önsöz metinizini yazınız.
\input{./front/prfc.tex}
%----------------------------------------------------------------------------------------------------
%----------------------------------------------------------------------------------------------------

\tableofcontents
\addcontentsline{toc}{section}{\footnotesize \scshape İç\.\i ndek\.\i ler}

\newpage
%----------------------------------------------------------------------------------------------------
%----------------------------------------------------------------------------------------------------
% 		Lütfen "ozetabs.text" dosyasına girerek bilimsel çalışmanızın özetini ve abstract 
%									metinizini yazınız. 
\input{./front/abstrct.tex}
%----------------------------------------------------------------------------------------------------
%----------------------------------------------------------------------------------------------------

\listoffigures
\addcontentsline{toc}{section}{\footnotesize \scshape Şek\.\i ller L\.\i stes\.\i}
\newpage
\listoftables
\addcontentsline{toc}{section}{\footnotesize \scshape Tablolar L\.\i stes\.\i}
\newpage


%----------------------------------------------------------------------------------------------------
%----------------------------------------------------------------------------------------------------
%	Bilimsel metninizde ekler listesi ve/veya simgeler listesi varsa lütfen "eklerlistesi.tex" ve
%	"simkislist.tex" dosyalarına girerek kendi listenizi yazınız. Her iki listede tablo biçiminde
%	oluşturulmuştur. Sayfa yapısı ve tablo ayarlarını değiştirmeden, tablo içerine kendi metin 
%	bilgilerinizi girerek size ait listeleri rahat bir biçimde oluşturabilirsiniz. 
%
%	Simge ve kısaltmalar isteniz 1 sayfadan fazla ise "ana.tex" dosyasında tanımlanmış olan "longtable" 
%	pakekini kullanmanızı tavsiye ederiz. Eğer kendinize ait yeni bir simge ve kısaltmalar listesi 
%	oluşturmak istiyorsanız enstitü tarafından yayınlanan TEZ YAZIM KLAVUZUNA UYMANIZ GEREKMEKTEDİR!
%	
%	Ekler listesi  manuel olarak tanımlandığından dolayı bu listeye  kendiniz doldurmalınız. Listede
%	sadece  ek numarası, ek ismi ve ekin sayfa numarısı bilgilerini doldurmanız yeterli olacaktır.
%	Yazım sitili ve boşluklar üzerinde herhangi bir değişiklik yapmanıza gerek bulunmamaktadır.
\input{./front/apndx.tex}
\input{./front/symabvr.tex} 
%----------------------------------------------------------------------------------------------------
%----------------------------------------------------------------------------------------------------



%----------------------------------------------------------------------------------------------------
%----------------------------------------------------------------------------------------------------
%		Bölüm, alt bölüm, 2. derece alt bölüm başlıkları format ve stil ayarları
%----------------------------------------------------------------------------------------------------
%----------------------------------------------------------------------------------------------------

\titleformat{\chapter}[hang]{\Large\bfseries\scshape}{\thechapter}{0.4cm}{\Large\bfseries\scshape}
\titlespacing{\chapter}{0pt}{-1cm}{8mm}

\titleformat{\section}[hang]{\normalsize\bfseries}{\thesection}{0.25cm}{\normalsize\bfseries}
\titlespacing{\section}{0pt}{\parskip}{4mm}
\titleformat{\subsection}[hang]{\small\bfseries}{\thesubsection}{0.2cm}{\small\bfseries}
\titlespacing{\subsection}{0pt}{\parskip}{4mm}
%----------------------------------------------------------------------------------------------------
%----------------------------------------------------------------------------------------------------

\pagenumbering{arabic}

%----------------------------------------------------------------------------------------------------
%----------------------------------------------------------------------------------------------------
%	Bilimsel metninize ait ana bölümlerin "ana.tex" dosyasına eklendiği bölüm. Aşağıda yer alan giriş 
%	dosyalarını kendi bilimsel metnize göre değiştiriniz. Aşağıdaki giriş dosyalarını kendi ihtiyacınıza
%	göre arttırıp azaltabilirsiniz.
%			
%
%	body klasörü içerisindeki dosyaların adlarını kendinize göre düzenlerseniz, lütfen aşağıdaki dosya isimlerini de onlarla aynı olacak şekilde güncelleyiniz.  
\input{./body/intro.tex} %Giriş
\input{./body/bolum2.tex} %Takip eden bölüm 
\input{./body/bolum3.tex} %Takip eden bölüm
\input{./body/bolum4.tex} %Takip eden bölüm
\input{./body/snclr.tex} %Bulgular
\input{./body/onrlr.tex} %öneriler
\input{./bib/bib.tex} %Bibliyografya DÜZENİ, lütfen DEĞİŞTİRMEYİNİZ.
\input{./back/appndx.tex} % Ekler listeniz yoksa lütfen bu komutu pasif pozisyona getiriniz!  
\input{./back/cv.tex} % CV Düzeni, lütfen değiştirmeyiniz.
%\printbibliography
\end{document}
